\documentclass[a4paper, 12pt]{article}

% --- PAQUETES ---
\usepackage[utf8]{inputenc}
\usepackage[spanish]{babel}
\usepackage{graphicx}
\usepackage{geometry}
\usepackage{hyperref}
\usepackage{titlesec}
\usepackage{parskip}
\usepackage{amsmath, amssymb}
\usepackage{float}
\usepackage{caption}
\usepackage{subcaption}
\usepackage{listings}

% Configuración de márgenes estándar
\geometry{
 left=2.5cm,
 right=2.5cm,
 top=2.5cm,
 bottom=2.5cm
}

\begin{document}

% --- CARÁTULA ---
\begin{titlepage}
    \centering

    \includegraphics[width=0.4\textwidth]{logo_udesa_negro.png}
    \par\vspace{1cm}

    {\scshape\Large Aprendizaje Reforzado \par}
    \vspace{0.5cm}
    {\large Docentes: Claudio Pose, Gabriel Gregorio Torre \par}
    \vspace{1.5cm}

    {\large\bfseries Trabajo Práctico Final \par}

    {\Huge\bfseries S.P.I.D.E.R. \par}
    \vspace{0.5cm}

    {\large\itshape System for Proximal Intelligence in Distributed Environments with ROS2 \par}
    \vspace{2.5cm}

    \begin{flushleft}
        \large
        \textbf{Autores:} \\
        Lisandro Morales Arce \\
        Facundo Gabriel Vulcano \\
        Ciro Russi \\
    \end{flushleft}

\end{titlepage}

% --- ÍNDICE ---
\newpage
\tableofcontents
\newpage

% --- CONTENIDO DEL INFORME ---

\section{Introducción}
En el presente informe se detalla el diseño, implementación y evaluación de un agente basado en aprendizaje reforzado para el control de un robot hexápodo en un espacio bidimensional, sobre el cual se encuentra un ``target'' que la araña tendrá como objetivo.\\
El trabajo consistió en entrenar un agente que simule el entorno sobre el que puede accionar la araña para alcanzar el target con el set de acciones óptimo, usando el algoritmo de PPO. Y luego, interconectar la política entrenada con los dispositivos de hardware requeridos para que la araña actúe en el espacio físico.

\section{Arquitectura del Sistema}
Esta sección describe los componentes físicos y lógicos que integran el proyecto S.P.I.D.E.R., incluyendo la comunicación inalámbrica y la estructura de nodos en ROS 2.

\begin{figure}[htbp]
    \centering
    \includegraphics[width=0.8\linewidth]{spider_comunication.png}
    \label{fig:myfigure}
\end{figure}

La araña se trata de un robot de seis patas que consta con un controlador que recibe comandos vía UART en un formato de bytes particular para el modelo. Dichos comandos solo pueden ser números enteros que se corresponden a un conjunto de acciones o ``motions'' que la araña puede llevar a cabo.\\\\
Para enviar dichos comandos en UART utilizamos una placa con microcontrolador modelo ESP32. Desde una notebook, se envían los números enteros asociados a los distintos comandos para accionar la araña hacia la ESP32 y ella es la que se encarga de traducir el entero en la tira de bytes con el formato esperado por el controlador de la araña, y luego enviar dicho paquete de datos.
\\\\
La notebook es la que se encarga de computar los movimientos de la política según las observaciones disponibles. Las observaciones se obtienen del un conjunto de cámaras operadas con el programa ``OptiTrack''; este programa permite triangular la posición de objetos en tiempo real. En el contexto de este proyecto, OptiTrack permitió obtener la posición de la araña y del objetivo en un espacio de coordenadas tridimensional. Dichas posiciones eran compartidas con la notebook para que se encargue de procesarlas y luego hacer la predicción de los movimientos.

\section{Entrenamiento de la Política}

El núcleo del comportamiento autónomo del agente reside en una política entrenada mediante el algoritmo \textit{Proximal Policy Optimization} (PPO). Para la implementación de este algoritmo utilizamos la librería \texttt{Stable-Baselines3} \cite{stable-baselines3}, que provee una versión robusta y ampliamente utilizada de PPO, así como herramientas estandarizadas para el entrenamiento y evaluación de agentes en entornos de tipo \texttt{Gym}.

En nuestro caso, el \textbf{espacio de observación} está definido por la posición relativa del target en la terna del robot. En la implementación final, la observación entregada a la red neuronal consiste en un vector continuo que codifica la distancia y dirección del objetivo respecto del agente (por ejemplo, las coordenadas $(x, y)$ del target en el marco local del robot y, en versiones extendidas, el ángulo de desalineación). Este vector se normaliza y se proporciona como entrada a la política de PPO, que lo procesa mediante una red neuronal \textit{feed-forward}.

El \textbf{espacio de acción} es discreto y corresponde directamente a los comandos de movimiento disponibles en la araña física. Se definieron 12 acciones, cada una asociada a una instrucción concreta. En el entorno de gym, este espacio se modela como un \texttt{Discrete(12)}, de manera que la política de PPO produce una distribución categórica sobre estas 12 acciones posibles en cada paso de tiempo.

Para optimizar el desempeño del agente, se llevó a cabo una \textbf{búsqueda de hiperparámetros} utilizando la infraestructura de \texttt{RL Zoo} (\texttt{rl-zoo3}) \cite{rl-zoo3}. Esta herramienta actúa como un envoltorio sobre \texttt{Stable-Baselines3} y permite automatizar experimentos de entrenamiento variando hiperparámetros clave de PPO, tales como la tasa de aprendizaje, el tamaño de lote (\textit{batch size}), el número de pasos por actualización (\textit{n\_steps}), el factor de descuento $\gamma$, el coeficiente de entropía y el rango de recorte (\textit{clip range}), entre otros. Mediante esta búsqueda sistemática se seleccionó una configuración de hiperparámetros que equilibrara estabilidad y velocidad de convergencia, mejorando la calidad de la política aprendida en el entorno propuesto.

\subsection{Definición del Entorno (MDP)}
Para llevar adelante el proyecto fue necesario desarrollar un entorno propio que capturara de manera adecuada la interacción entre el agente (la araña) y el objetivo (el target). El primer paso consistió en definir las dimensiones del mapa, tomando como referencia el espacio físico real donde posteriormente se realizarán las pruebas. Esto permitió establecer un área maniobrable consistente y reproducible dentro de la simulación.

El estado del entorno se modeló a partir de dos objetos principales: la araña y el target. Para simplificar el diseño y facilitar el aprendizaje, se fijó la posición de la araña en el origen del sistema de coordenadas locales. De este modo, la observación del agente se reduce a las coordenadas relativas del target respecto de la araña, expresadas como $(x,y)$ en la terna local del robot. Esta representación permite capturar de forma directa la relación espacial relevante para la tarea, evitando la necesidad de trabajar con posiciones absolutas en el mapa y reduciendo la complejidad del estado.

Las acciones disponibles para el agente están directamente vinculadas a los comandos reales que puede ejecutar la araña física. Cada movimiento corresponde a un código específico (por ejemplo, avanzar utiliza el comando \texttt{252}). Con el fin de mantener compatibilidad entre la simulación y el hardware final, se definió un espacio discreto de 12 acciones, donde cada una representa un comando de movimiento permitido por la plataforma.
Estas acciones pueden agruparse en los siguientes tipos de maniobras fundamentales:

\begin{itemize}
    \item \textbf{Rotaciones puras:} permiten modificar la orientación del robot sin producir desplazamiento lineal.
    \item \textbf{Movimientos hacia adelante y atrás:} incluyen variantes rectas y con giro suave.
    \item \textbf{Movimientos rápidos:} versiones de mayor velocidad para algunas de las acciones.
\end{itemize}

La función de recompensa fue diseñada para fomentar un comportamiento eficiente y estable durante la navegación hacia el objetivo. El agente recibe una combinación de señales que reflejan tanto su progreso espacial como la calidad de sus movimientos. La recompensa total se construye a partir de los siguientes cinco términos:

\begin{itemize}
    \item \textbf{Progreso en distancia:} se calcula la diferencia entre la distancia previa y la distancia actual al target. Acercarse genera un refuerzo positivo, mientras que alejarse produce un valor negativo. Este término se encuentra escalado por el tamaño del mapa para mantener una magnitud consistente.

    \item \textbf{Costo base por paso:} cada acción implica un pequeño castigo negativo. Este término desalienta que el agente prolongue el episodio sin necesidad y promueve la eficiencia temporal.

    \item \textbf{Mejora de orientación:} si el agente logra reducir la distancia, se evalúa si también mejoró su alineación con la posición del target. El cambio angular se normaliza y se pondera mediante un factor que depende de la distancia actual, de modo que la orientación sea más relevante cuando el objetivo está lejos.

    \item \textbf{Penalización por retroceso:} si el agente aumenta su distancia al objetivo, recibe un castigo adicional. Esto evita comportamientos erráticos, oscilaciones o trayectorias innecesariamente largas.

    \item \textbf{Bonus por éxito:} al alcanzar el radio de éxito definido en el entorno, el agente recibe una recompensa positiva significativa, incentivando completar la tarea de navegación.
\end{itemize}

\begin{equation}
    R_t =
    \underbrace{\Delta d_t \cdot \alpha}_{\text{progreso en distancia}}
    \;-\;
    \underbrace{c}_{\text{costo por paso}}
    \;+\;
    \underbrace{\mathbb{1}_{\Delta d_t > 0}\;\beta \;\Delta \theta_t \; s(d_t)}_{\text{mejora de orientación}}
    \;-\;
    \underbrace{\mathbb{1}_{\Delta d_t < 0}\;\gamma}_{\text{penalización por alejarse}}
    \;+\;
    \underbrace{\mathbb{1}_{d_t \leq r_{\text{succ}}}\;B}_{\text{bonus de éxito}} ,
\end{equation}

Esta combinación de términos permite que la política aprendida priorice una aproximación estable: orientarse correctamente, avanzar de forma eficiente, evitar retrocesos y completar episodios con comportamiento consistente.

\section{Desarrollo Experimental}
La transición del modelo simulado al hardware real presentó desafíos significativos, particularmente en la transformación de coordenadas entre el sistema global de referencia y el marco local del agente; tuvimos que hacer mediciones de movimiento para cada comando a mano, lo cual agregó incerteza al problema.

En cuanto al software, el primer aspecto abordado en el desarrollo del proyecto fue la construcción del entorno. Desde el inicio se decidió simplificar el espacio de estados con el objetivo de facilitar y acelerar el entrenamiento. Para la tarea de navegación considerada, se determinó que la información esencial era la \textbf{pose relativa del target respecto del robot}, compuesta por las coordenadas $x$, $y$ y el ángulo $\theta$.

Bajo esta representación, el robot se asume siempre en el origen del sistema de coordenadas local, mientras que la observación corresponde exclusivamente a la posición del target transformada desde la terna global a la terna del robot. De este modo, el agente percibe únicamente la distancia y dirección relativas al objetivo, eliminando información redundante y reduciendo la dimensionalidad del espacio de estados.

Una vez definida la estructura de la observación, el siguiente paso fue diseñar una función de recompensa adecuada. La primera aproximación consistió en utilizar la distancia negativa al target junto con un término que premiara la alineación del robot hacia el objetivo. No obstante, esta combinación presentó un problema: el agente encontró una manera de ``hackear'' la recompensa. En lugar de completar la tarea, aprendía a acercarse al target y luego realizar pequeñas oscilaciones angulares, aprovechando el término de orientación para acumular recompensa sin llegar nunca a la región de éxito.
\begin{figure}[H]
    \begin{subfigure}[b]{0.5\textwidth}
        \centering
        \includegraphics[width=1\textwidth]{curva_reward.jpg}
        \caption{Reward por episodio}
    \end{subfigure}\hfill
    \begin{subfigure}[b]{0.4\textwidth}
        \centering
        \includegraphics[width=1\textwidth]{curva_steps.jpg}
        \caption{Cantidad de pasos por episodio}
    \end{subfigure}
    {\small
    \caption{Steps degenerados y recompensa constantemente mejorada.}
    \label{fig:comparacion_alg}}
\end{figure}
Posteriormente reformulamos la función de recompensa para evitar los problemas detectados. En lugar de utilizar valores absolutos de distancia y orientación, la nueva versión emplea la \textbf{diferencia entre el paso previo y el paso actual}, tanto en distancia como en alineación. De este modo, el agente es recompensado únicamente por mejorar su situación relativa respecto del target, eliminando la posibilidad de explotar el término de orientación sin avanzar hacia el objetivo.

Si bien esta modificación resolvió el problema inicial, reveló un nuevo comportamiento indeseado durante el análisis del \textit{policy map}. El agente comenzaba a elegir sistemáticamente la acción de desplazarse rápidamente hacia atrás. Observamos que el robot solo realizaba rotaciones y luego ejecutaba la acción de \emph{fast backward} de manera repetida. La causa estaba en un error de configuración: habíamos establecido por equivocación un desplazamiento de 70cm para esa acción, en lugar de los 7cm reales. Debido a esta discrepancia, el agente evaluaba esa acción como altamente beneficiosa, ya que producía un cambio de posición considerable dentro de la simulación.

\begin{figure}[H]
    \centering
    \includegraphics[width=1\textwidth]{policy_map.png}
    \caption{Numero 9 backward, 0 pivot left y 1 pivot right}
\end{figure}

Una vez corregido este valor y mejoradas las mediciones reales mediante el sistema de captura de movimiento OptiTrack, la política aprendida dejó de mostrar este sesgo y comenzó a utilizar el conjunto completo de acciones de forma coherente con la dinámica física del robot.

Finalmente, antes de obtener la política definitiva, detectamos un error en la configuración de los comandos angulares: las acciones de \textit{pivot right} y \textit{pivot left} estaban intercambiadas en la implementación. Esto implicaba que, al seleccionar una acción de giro hacia la derecha, el robot en realidad rotaba hacia la izquierda, y viceversa. Como consecuencia, parte del comportamiento aprendido por la política estaba adaptado a una dinámica de rotación incorrecta. Una vez identificado el problema, corregimos la asignación de estos comandos y fue necesario reentrenar el agente, incorporando estas dificultades como parte del proceso iterativo de depuración entre el entorno simulado y el control real del robot.

\section{Análisis de los Resultados}
Se analizó el desempeño final a partir de tres fuentes: las corridas de \textit{RL Zoo}, los mapas de política obtenidos durante el entrenamiento y las pruebas de transferencia a hardware.

\textbf{Curvas de entrenamiento.} En la búsqueda de hiperparámetros con \textit{TPE median pruner} se observaron dos comportamientos claros en los reportes \texttt{logs/ppo/report\_*.csv}: (1) varias combinaciones divergían o quedaban estancadas en recompensas muy negativas, y (2) un subconjunto de configuraciones con $\gamma$ alto y lotes pequeños convergía de forma estable. La mejor corrida completada alcanzó una recompensa promedio de episodio positiva (35.18) manteniendo el número de pasos acotado, lo que indicó que el shaping diferencial de distancia y orientación resolvió el problema de oscilaciones sin progreso. Sobre esa región del espacio de hiperparámetros se fijó la configuración definitiva usada en \texttt{spider.py}.

\textbf{Mapas de política.} Las figuras generadas automáticamente (ej. \texttt{policy\_map.png}) mostraron la estrategia espacial aprendida. En la versión final, el mapa se organiza de manera interpretable: el frente del robot concentra acciones de avance rápido y recto, mientras que los sectores laterales activan pivotes y combinaciones de giro suave. Tras corregir la escala del comando de \emph{fast backward}, la política dejó de presentar zonas extensas asociadas a retroceso y distribuyó el uso de las 12 acciones, lo que sugiere una navegación más balanceada. También se redujeron las discontinuidades abruptas entre regiones de acción, señal de mayor suavidad en la política.

\textbf{Transferencia \textit{sim-to-real}.} Al conectar la política entrenada con el pipeline OptiTrack--ROS2--ESP32, el comportamiento observado en hardware siguió las tendencias del simulador: el robot se orienta primero hacia el objetivo y luego alterna entre avances rectos y giros cortos. Los fallos principales provinieron de desalineaciones de terna y del retardo entre lectura de posición y envío de comando; ambos efectos introdujeron una deriva que a veces desviaba la trayectoria cerca del objetivo. El uso de recompensas diferenciales ayudó a corregir estos desvíos en línea, pero persiste una brecha cuando hay ruido en las mediciones o superficies con poca fricción.

\textbf{Obstáculos y colisiones.} La versión del entorno utilizada para los entrenamientos finales incluye obstáculos circulares aleatorios, reportados en la observación como las posiciones relativas de los cinco más cercanos. Durante las primeras fases, las colisiones eran frecuentes; al avanzar el entrenamiento la política aprendida tiende a rotar y buscar rutas laterales antes de avanzar rápido, lo que se refleja en menos terminaciones por choque. En hardware no se reprodujeron obstáculos físicos, por lo que esta capacidad quedó validada sólo en simulación y se identifica como un punto a completar en trabajos futuros.

\textbf{Calibración y sensibilidad.} El sistema resultó sensible a pequeñas inconsistencias en los parámetros de movimiento: la inversión de pivotes y el error de escala en el retroceso rápido cambiaron radicalmente el mapa de acciones. Esto resalta la necesidad de contar con mediciones confiables de desplazamientos y rotaciones (OptiTrack ayudó a reducir la varianza) y de mantener sincronía temporal entre el entorno simulado y el real. Aun con estas limitaciones, la política final mostró estabilidad frente a ruido moderado en las observaciones y generalizó sobre distintas posiciones iniciales del target dentro del mapa.

\section{Conclusiones}

El proyecto \textsc{S.P.I.D.E.R.} demostró de manera consistente la viabilidad de controlar un robot de elevada complejidad mecánica mediante políticas de Aprendizaje por Refuerzo entrenadas íntegramente en simulación y transferidas con éxito al hardware real. El desarrollo de un entorno personalizado, la definición de un espacio de estados minimalista y la implementación del algoritmo PPO utilizando \texttt{Stable-Baselines3} permitieron obtener comportamientos estables y reproducibles bajo condiciones físicas reales.

A lo largo del proceso debimos superar diversos desafíos, entre ellos la transformación precisa entre la terna global del sistema OptiTrack y la terna local del robot, la calibración de los comandos motores y la corrección de inconsistencias en la dinámica de rotación y traslación. La incorporación de mediciones fiables mediante OptiTrack y la búsqueda de hiperparámetros con \texttt{RL Zoo} resultaron fundamentales para mejorar la estabilidad del entrenamiento y reducir la brecha entre simulación y realidad.

Si bien el sistema funcionó correctamente, el proyecto presenta ciertas limitaciones: el simulador no modela con fidelidad la física real del robot, el entorno carece de obstáculos y no se midió el \textit{wall time} asociado a cada acción, lo que genera discrepancias en la temporalidad real del movimiento. Aun así, los resultados muestran que, incluso con un estado extremadamente simplificado y un conjunto reducido de acciones, el agente es capaz de navegar eficazmente hacia un objetivo utilizando únicamente información relativa local. Esto demuestra que políticas de RL bien diseñadas pueden controlar robots reales pese a restricciones sensoriales y mecánicas.

Más allá del éxito técnico, el proyecto constituye una contribución valiosa para quienes deseen continuar investigando la interacción entre aprendizaje por refuerzo y robótica física. La metodología desarrollada que integra simulación, calibración mediante captura de movimiento, transferencia \textit{sim-to-real} y entrenamiento offline configura un marco replicable y extensible para estudios futuros.

En conjunto, \textsc{S.P.I.D.E.R.} cumple sus objetivos y abre múltiples líneas de trabajo, como la incorporación de entornos más complejos, sensores adicionales, políticas multiobjetivo, estrategias más robustas al ruido y la exploración de control multirobot. El proyecto no solo valida un enfoque, sino que establece una base sólida para futuras investigaciones en navegación autónoma mediante Aprendizaje por Reforzamiento.

\clearpage
\bibliographystyle{plain}
\bibliography{biblio}
\end{document}
